\section{Experimental Setup}


\subsection{Wearable Thermal Data Collection Platform}

%Paragraph 1 — Overview and Motivation
%Introduce the device and design goals.
We developed a wearable thermal sensing platform designed for privacy-preserving indoor navigation research. The platform consists of multiple ultra-low-resolution thermal cameras mounted on a wearable frame and connected to an embedded computing unit for synchronized data acquisition. The design emphasizes low cost, privacy preservation, and flexibility, making it suitable for studying navigation perception under severe sensing constraints.


%Paragraph 2 — Hardware Description
%Describe the hardware without excessive detail.

The wearable platform incorporates multiple MLX90640 SparkFun IR Array Breakout sensors, each providing a $32\times24$ thermal image of the environment. The sensors are mounted at different viewing angles to increase environmental coverage and capture complementary observations of the user's surroundings. All sensors are connected to an NVIDIA Jetson Nano-based acquisition unit that records timestamped thermal frames for subsequent processing.

% Paragraph 3 — Data Collection Procedure
% Explain how data are acquired.



\subsection{Data Collection}

Data were recorded by walking the wearable rig through indoor environments while
a high-resolution RGB camera, rigidly mounted alongside the thermal array,
captured a time-synchronized reference stream used only for label generation.
During each traversal the three active thermal cameras and the RGB camera
recorded timestamped frames continuously, producing sequences that span clear
corridors as well as encounters with people and static objects, so that all
three navigability states are represented across the left, middle, and right
directions. Only the thermal tiles and their automatically generated
per-direction labels (Section~\ref{sec:labels}) are retained for model training;
the RGB stream is discarded after labeling and is never available at inference.

\subsection{Training Protocol}

To avoid leaking temporally adjacent, near-duplicate frames across splits, we use
a blocked partition: each recording is divided into contiguous blocks and the
final $20\%$ of every block is held out for validation, giving a train/validation
split that keeps correlated neighbors together. Camera-identity-shuffle
experiments are additionally evaluated on a held-out test set. All learned models
share the same masked, per-camera, three-class cross-entropy objective and are
trained with Adam (learning rate $10^{-3}$, weight decay $10^{-4}$) under a
cosine-annealing schedule for $30$ epochs at batch size $64$, with class
weighting to counter the imbalance toward clear directions. For Feature Fusion we
additionally sweep backbone depth, width, activation, and head capacity and
report the best configuration. We evaluate with per-class $F_1$ score reported
per direction (left, middle, right) and overall; because the label distribution
is skewed toward clear directions, macro-averaged $F_1$ is used for model
selection rather than raw accuracy.